\chapter{Background}
\label{ch:background}

\section{RFI in spectrograms}
\label{sec:rfi_types}
In a spectrogram, different kinds of RFI leave different shapes, depending on
how long they last and how many frequencies they cover.
\begin{itemize}
  \item \textbf{Constant narrowband RFI} affects a single frequency channel for
    the whole observation and appears as a thin line along the time axis.
    Typical sources are fixed terrestrial transmitters, such as a particular
    radio or television station.
  \item \textbf{Constant broadband RFI} covers a wide range of frequencies for the
    whole observation and appears as a thick band. Continuously broadcasting
    satellites are a common cause.
  \item \textbf{Instantaneous wideband RFI} lasts only a moment but covers almost
    all frequencies, appearing as a thin line along the frequency axis.
    Lightning, arcing on power lines and radar sweeps produce this kind.
  \item \textbf{Patchy broadband RFI} covers several channels for a few minutes
    and then fades, appearing as irregular, cloud-like patches.
\end{itemize}
Not every structure in a spectrogram is RFI. The receiver's sensitivity
varies with frequency (the bandpass, \cref{sec:bandpass}), and slow drifts and
standing waves add smooth patterns. A flagging method has to distinguish these
from real interference.

\section{Classical flagging methods}

\subsection{Statistical thresholds}
The simplest approach flags any pixel more than $k$ standard deviations above
the mean. Tools used for folded pulsar data, such as \texttt{paz} and
\texttt{clfd}, are based on robust statistics of this kind, such as medians.
They need no training, but they assume the clean data is close to Gaussian, so
they miss faint RFI and struggle when the noise level changes across the data.
On the LOFAR test data, the best simple threshold of this kind reaches an
\Fone{} of 0.41 (\cref{sec:tfunet_lofar}).

\subsection{AOFlagger}
\aoflagger{}~\cite{offringa2012} is the standard RFI flagger for LOFAR. It
combines thresholding at several scales with a morphological step that
extends flags along time and frequency, so that the edges of RFI features are
also caught. It is important for this report because it produced the
training labels of the LOFAR dataset (\cref{sec:lofar}).

\subsection{The 2\texorpdfstring{$\sigma$}{-sigma}CRF method}
Chen et al.~\cite{chen2023} clean RFI in folded pulsar data in two steps.
First, they compute the root-mean-square of the folded profile in each
frequency channel and time sub-integration, giving a 2D image, remove its
baseline, and fit a Gaussian to the histogram of its values. Pixels within
$2\sigma$ of the Gaussian peak are labelled clean, pixels beyond $3\sigma$ are
labelled RFI, and those in between are left ambiguous. Second, they refine
these labels with a fully connected conditional random field
(CRF)~\cite{krahenbuhl2011}. A CRF treats every pixel as connected to the
others and finds the labelling that minimises an energy made of two parts: a
unary term, how likely each pixel is to be RFI on its own, and a pairwise
term, which encourages pixels that are close together and similar in value to
share a label. This lets weak RFI that is unrecognisable in a single pixel be
picked out from its neighbourhood. They report 93--95\% accuracy on data from
the FAST telescope. The method still rests on a Gaussian model of the clean
data, which is where a learned approach might do better.

\section{Neural networks for RFI segmentation}

\subsection{U-Net}
The U-Net~\cite{ronneberger2015} is a convolutional network for image
segmentation. Its encoder repeatedly applies convolutions and halves the image
size, building up a description of \emph{what} is in the image; its decoder
reverses this, restoring the full size to locate \emph{where} it is. Skip
connections pass fine detail directly from each encoder level to the matching
decoder level. The output has one probability per pixel---here, the
probability that the pixel is affected by RFI.

\subsection{\tfunet{}}
Akeret et al.~\cite{akeret2017} applied a U-Net to RFI in their \tfunet{}
code, written in TensorFlow. The configuration used in this report has 3
levels and 32 filters in the first layer, giving 465{,}986 trainable
parameters. Two design details become important later: \tfunet{} applies a ReLU
to its output scores, and it has no normalisation layers
(\cref{ch:tfunet}).

\subsection{Recent methods}
Mesarcik et al.~\cite{mesarcik2022}, who released the LOFAR dataset used here,
proposed an approach called Nearest-Latent-Neighbours, which learns only from
RFI-free data and flags whatever looks unusual. They compared it with several
supervised networks on the same LOFAR test set. The best \Fone{} in their
comparison came from RFI-Net~\cite{yang2020rfinet}, a residual U-Net, with
0.5979; their own U-Net reached $0.5876 \pm 0.0031$. Their method had the best
area under the ROC curve but a lower \Fone{} of 0.5114. These published values
are the benchmarks used in \cref{ch:hybrid}.

\section{Measuring performance}
\label{sec:metrics}
For each pixel, a flagging method either flags it as RFI or not, and it either
is or is not RFI in the reference labels. Counting the true positives (TP),
false positives (FP) and false negatives (FN) over all test pixels gives
\begin{equation}
  \text{precision} = \frac{TP}{TP + FP}, \qquad
  \text{recall} = \frac{TP}{TP + FN}, \qquad
  F_1 = \frac{2\,\text{precision}\cdot\text{recall}}{\text{precision} + \text{recall}}.
\end{equation}
Precision is the fraction of flagged pixels that really are RFI; recall is the
fraction of RFI that was found. \Fone{} is high only if both are.

Accuracy, the fraction of pixels classified correctly, is not useful here.
RFI covers less than 1\% of the LOFAR test pixels, so a method that flags
nothing at all is 99.23\% accurate while finding no RFI.

The area under the ROC curve measures something different: how well a model
\emph{ranks} RFI pixels above clean ones, regardless of where the final cut is
made. A model can rank well and still make poor final decisions
(\cref{sec:hybrid_lofar}).

\paragraph{Choosing the threshold.} A network outputs a probability for each
pixel, and a threshold turns it into a yes/no decision. The choice of threshold
changes the \Fone{}, and there are two ways to make it.
\begin{itemize}
  \item \textbf{Maximum \Fone{}:} try every threshold on the test set and report
    the best. This is optimistic, because it uses the test labels to pick the
    threshold. It is, however, what Mesarcik et al.\ report, so it is needed
    for a fair comparison with their results.
  \item \textbf{Pooled \Fone{}:} choose the threshold on a separate validation
    set, fix it, and then apply it to the test set. This is the honest
    measure of how a model would perform on new data.
\end{itemize}
For example, if a model scores ten pixels and the best threshold on those ten
gives $F_1 = 0.89$, but the threshold chosen beforehand on validation data
gives $0.86$ on the same pixels, the maximum \Fone{} is 0.89 and the pooled
\Fone{} 0.86. The maximum is never lower than the pooled value. In this report
both are given; the pooled value is the one I rely on.

``Pooled'' also means that TP, FP and FN are added up over all test images
before computing a single \Fone{}. Averaging a separate \Fone{} for each image
gives a different number: for \aoflagger{} on LOFAR, the per-image average is
0.103 lower than the pooled value.
